\section{Measurement and the registered claim}\label{sec:measurement}

This section states how the study measures and tests its claim, as the simulation code computes it and as the
project's record of decisions fixed it before any confirmation run. The contract is in Section~\ref{sec:contract} and
the model in Section~\ref{sec:model}; nothing here changes either. Every quantity is a paired difference between
two arrangements run on the same simulated people and the same random draws. One agent is one consumer unit, so
money is reported in dollars a person a year unless stated.

\subsection{The experiment}

Three controls vary over a full cross of 105 cells:
\begin{enumerate}[nosep]
  \item the reserve rate $\rho$, equal to the outside rate $r_{\mathrm{out}}$: 2, 4.4 and 6\% a year;
  \item launch participation $a_{0}$, the share of people in the network at launch: 0.05, 0.1, 0.2, 0.3, 0.45,
        0.6 and 0.8;
  \item the intensity of a distress year, the extra monthly rate at which jobs lose their wage in months 25 to 36
        after launch, above normal job loss of 0.55\% of jobs a month: 0, 0.5, 1, 1.5 and 2 points.
\end{enumerate}
The five distress levels are named by these points, 0 to 2. At month 36 the share of jobs without a wage is
about 1.9, 3.4, 4.9, 6.3 and 7.7\% in that order. The rule for the wage of an extra layoff is in
Section~\ref{sec:model} and is the same at every level. The 2-point level is the headline of the stress clause and
is labelled 2009's wage loss before unemployment insurance.

The whole grid, run on a set of seeds, is called a map. A run is one cell on one seed, in five arrangements:
the bank arrangement (bank money, cards and outside lenders); best-case DeFi (stablecoin payments and a
reserve-funded yield, and no credit for a payer who is short, Section~\ref{sec:model}); the pool without the charge
(the payment money $P$ and the household-owned pool, with no demurrage); CPM; and indexed CPM, in which $P$ is
indexed to consumer prices. Indexed CPM is defined only at a positive real reserve rate, so its cells at
$\rho = 2\%$, where the real rate is zero at 2\% inflation, are its declared limit and stay outside its region.

The arrangements share common random numbers and the burn-in (Section~\ref{sec:model}), so wealth at launch, and
with it every group used below, is the same in all five.
Within a cell and seed the runs at different distress levels are identical up to the start of the distress year,
and the bank arrangement does not read $a_{0}$. Both facts are checked bit for bit against direct runs before
every map, and the simulation code refuses to proceed otherwise.

Of the 60 months after launch (Section~\ref{sec:model}) two windows are read: the first covers months 13 to 36
and the second months 37 to 60. The rest of the paper calls the second window years 4 and 5. Gains and real cost
are summed over all 60 months after launch and expressed a year. The contract's state, participation and owners' return are read by
window, and where one figure is given it is the second window's mean.

\subsection{Outcomes}

Table~\ref{tab:outcomes} defines each measured quantity. A difference is always the named arrangement less its
yardstick, the bank arrangement or best-case DeFi, so a positive gain or saving favours the arrangement
named and a negative difference in failed payments favours it.

\begin{table}[ht]
\centering\small
\caption{Measured quantities, as the simulation code computes them.}
\label{tab:outcomes}
\begin{tabularx}{\textwidth}{@{}l X l@{}}
\toprule
Quantity & Definition & Unit \\
\midrule
Failed share & Of the value of the bills and purchases that fell due in the second window, the part not paid on
  its due day. Wages are left out of both; a wage not paid on its due date stays owed and is recorded apart. The
  text calls it the share of payments that fail. & \% of value due \\
Failed payments & The unpaid share of each failed bill or wage, each obligation counted once whatever the fee or
  the monthly fee cap; the arrangement less the yardstick. & shares a person a year \\
Wealth gain & End net wealth after 60 months less the yardstick's, person by person, expressed a year and averaged
  over people. & \$ a person a year \\
Complete gain & The wealth gain plus the difference in purchases received (a person's own service and shock
  purchases completed as buyer and the price of each durable bought that settled), with sellers' extra sales counted
  at face value; at the zero end, the same less the difference in sales as seller. At a surplus share between the
  ends it is the zero end plus that share of the difference in sales. & \$ a person a year \\
Real cost & The resources an arrangement's payments and credit use, each service at one sourced rate whoever
  provides it: card-type credit 3.4\% and term credit 2.5\% of balances a year, \$2 for each failed or pool-rescued
  share of a bill or wage, the card rails' costs, chain fees, 0.24\% of balances held and \$2 an earned-wage
  advance. Reported as the saving, the yardstick's real cost less the arrangement's. & \$ a person a year \\
Owners' return less $a$ & Households' realised return as pool owners, the income credited to
  their shares over their capital-time, before inflation, less their alternative $a$, the reserve rate on the
  arrangement's own unit floored at zero plus a declared spread of 0.25 point. & \% a year, by window \\
Bridges served & The bridges the pool drew over the bridges payers asked of it. & share, second window \\
Holding charge $c$ & The demurrage rate on wallet balances of $P$, which funds the liquidity the pool keeps ready
  (Section~\ref{sec:contract}). & \% a year, second window \\
$\theta$ & The coupon reduction the charge paid on carried pool bridges, over that reduction plus the coupons
  borrowers accrued on them: demurrage's share of bridging finance. & share, second window \\
\bottomrule
\end{tabularx}
\end{table}

Owners' return against their alternative is reported with its seed mean and 95\% interval by window, and no verdict
reads it. Beside the real cost two read-outs are stored from counts already recorded. The break-even cost of one
rescued payment is the cost at which CPM and best-case DeFi would use the same real resources; it is computed as \$2 less
the difference in real cost over the difference in rescued shares, on the cell's seed means. The other read-out is the real cost per
dollar of purchases completed in each arrangement.

\subsection{The claim}

The claim is read for CPM and for indexed CPM separately, in each $(\rho, a_{0})$ cell at distress level 0, on the
seeds' paired differences. Every interval check uses two intervals over the $n$ seeds: the two-sided $t$ interval
\begin{equation}
  \bar{d} \pm t_{1-\alpha/2,\,n-1}\,\frac{s}{\sqrt{n}},
  \label{eq:tinterval}
\end{equation}
with $\bar{d}$ the seed mean and $s$ the spread of the paired difference across seeds, and the bootstrap percentile
interval of the seed mean over 10,000 resamples of the seeds on a fixed random stream. A check reads supported
when both intervals lie wholly on the clause's side of its margin, falsified when both lie wholly on the wrong
side, and inconclusive otherwise. It holds only when it reads supported. Each check is read at $\alpha = 0.05$ and at
the arrangement's Bonferroni level, 0.05 over its 21 cells for CPM and over its 14 valid cells for indexed CPM.

The clauses are evaluated as follows.
\begin{enumerate}[label=(\alph*), leftmargin=*, itemsep=2pt]
  \item \emph{CPM operates and pays its own way.} A seed counts when households hold at least 90\% of pool capital,
        net of exit requests still queued, in every month of the second window, and $P$ payments complete in every
        month of it. The clause holds when the Wilson lower bound of the share of such seeds is above one half, no
        seed records that the issuer would leave, no seed has issuer top-ups to the reserves, and the cell reads
        Persistent (below). The pool's own margin and the joint issuer-and-pool test, as the claim first worded them,
        are stored beside; with break-even coupons and free entry both centre on zero by construction.
  \item \emph{Real cost.} The lower bound of the real resources saved is at or above $-\$2$ a person a year against
        best-case DeFi and against the bank arrangement. A real saving is claimed only where both lower $t$ bounds
        are above zero.
  \item \emph{Completed service.} Failed payments are lower, and purchases completed higher, than under best-case
        DeFi. With no distress this holds in every cell, as it must once a pool lends without collateral, so it
        is reported and is not evidence. Under distress it can fail where owners run on the pool
        (Section~\ref{sec:stress}).
  \item \emph{No group loses.} Eleven groups are read: the four classes of card holding crossed with wage or
        business income, the five quintiles of wealth at launch and the short and long billing cycles. For every
        group the lower bound of its mean complete gain must be at or above $-\$2$ against the bank arrangement at
        both ends of the bracket and against best-case DeFi at the face end, 33 checks in all. A group defined in
        fewer than two seeds has no interval and fails the clause.
  \item \emph{Stress.} For a cell that passes (a), (b) and (d) at level 0, the levels 0.5, 1, 1.5 and 2 points are
        tested in that order at level 0's Bonferroni level, stopping at the first failure. A level is withstood when
        clauses (a) to (d) hold on its run and the upper bound of the dip is at most one failed payment a year
        against the bank arrangement's dip and against best-case DeFi's. The dip is each seed's cumulative excess of
        failed shares over months 25 to 60 over the same seed's run at level 0, divided by three. The clause reports
        the highest level withstood and whether the cell withstands 2 points; it is not part of passing.
\end{enumerate}

Against best-case DeFi the registered clause values completed purchases at their price: no class, wealth quintile
or billing-cycle group may be shown behind by more than the \$2 margin. On the confirmation map every group
is ahead on this reading in every cell (Section~\ref{sec:results}). The money-only reading, with completed purchases valued at zero, is negative
only for people who are never short, by the yield the holding charge takes, the premium for the pooled
buffer. Its break-even value is the smallest surplus share, on a grid of 0.01, from which every group's bound holds
at every larger share. That value is computed on the bounds and on the group means and decides no verdict; it is
stored with the results and is not reported here.

A cell passes when (a), (b) and (d) hold at 95\%, with (d) holding both as built and under the cash-advance bound
of the next subsection. A falsifier overrides a pass: CPM is read as worse than best-case DeFi when, in at least
80\% of seeds, its failed share in the second window is higher, or the median wealth of a group of launch-wealth
quintile and billing cycle falls short of best-case DeFi's by more than 2\% of that group's median wealth at
launch. It is also read as worse when, on both intervals, the interest rate borrowers pay over the principal they borrow is higher, or
owners' second-window return is below DeFi's savings rate. Together these four are the falsifier
``worse than best-case DeFi'', one of the six listed under persistence below, and any one of them falsifies the
cell. Under the first two, the seeds in which CPM is worse stop counting as Persistent; under the last two, which
are read across seeds, every seed does. The cell then fails (a) and stays outside the region. A pass is worded
``non-inferior real cost with more completed service''. It is conditional on three things: the pool, the $P$ rail
and the issuer bear the sourced cost of what they replace; the card fee is today's 1.7\%; and best-case DeFi lends
only against collateral.

The region is read in five slices, one per distress level, with level 0 the headline. At level 0 it is the largest
set of passing cells joined edge to edge in the $(\rho, a_{0})$ grid, corners not joining and ties named. A cell
the seed count leaves underpowered cannot enter it. At a stressed level a cell enters when its level-0 cell passes,
its highest level withstood reaches that level, and its stress leg is not underpowered. On a stressed run clause
(d) is read as built; the cash-advance bound enters a stressed slice only through the level-0 cell, which must
pass under it. ``Holds somewhere'' is the
set of level-0 cells that pass at the Bonferroni level.

Persistence is read seed by seed from the second window against the first. Use holds when participation is at
least the first window's less 0.05 and at least the launch share less 0.05. Lending recurs when the pool makes new
loans in at least 80\% of months and lends at least 80\% as much per participant. Owners stay when capital is at
least 80\% of the first window's and at least 95\% of exit value is paid within two months. The pool pays its own
way when it runs without outside money. Credit quality holds when the write-off rate and the arrears share each
rise by at most one point. A seed is Persistent when all five hold and no falsifier fires: collapse, credit death,
dependence on subsidy, worse than best-case DeFi, unsafe backing or write-offs, or an issuer that would leave. A cell
reads Persistent when the Wilson lower bound of its share of Persistent seeds is above one half. A second view,
stored beside the claim and read by no clause, combines persistence with the group test on $t$ bounds alone.

\subsection{Bounds reported beside the claim}

\emph{The cash-advance bound.} Every gain is also computed with the cash-advance fees and cash-advance interest the
person saves against the yardstick taken out, to first order with no change in behaviour, for everyone and for each
group. The bound exists because card holders' use of cash advances in the bank arrangement runs about 15 times the
United States level. As built, that excess leans toward CPM, because the fees and interest saved on those advances
count as gain; under the bound they do not. A cell counts in the region only when it passes both as built and
under the bound.

\emph{Money paid to outside lenders and intermediaries.} The efficiency claim rests on real cost alone, which reads
no destination of any flow; interest and fees are transfers. The household gain is reported three ways. As built,
outside lenders' income leaves the households, and every verdict reads this case; it leans toward CPM on wealth and
leaves real cost unchanged. At a central share of 0.10 of that income returning to households as owners of the
lenders (range 0.06 to 0.14), the return is added as a ledger line by wealth group, for means only and for no
verdict. With all of it returned, in proportion to spending, the outer bound also hands back the lenders' staff,
funding and loss costs, and so leans against CPM.

\emph{The loan-loss bound.} As built a written-off debt stays owed and is later collected, which leans toward CPM
on the bridge price. Under the bound every lender, the pool included, loses each month the published charge-off
rate of its loan class on its average balance, moving with the share of jobs without a wage three months earlier,
and the pool's coupons carry that rate. The bound leans against CPM.

\emph{Interest returned.} The outer bound above is run as its own map, with outside lenders' interest, and the
card's cash-advance and late fees, returned to households in proportion to spending. The loan-loss bound and this
one run over the whole map on the three development seeds of the final development map, and are reported as bounds
from development seeds.

\subsection{Sample sizes}

\emph{Population.} The number of agents is a numerical setting, fixed by a finite-size check whose rule was written
before the check was read. Populations of 500, 1,000, 2,000, 4,000 and 8,000 people were run on eight seeds kept for
nothing else, in three cells ($\rho = 4.4\%$; $a_{0} = 0.05$, 0.2 and 0.6) with every arrangement. The rule takes
the smallest size at which doubling the population moves none of twelve headline numbers in any of the three cells
by more than that number's own 95\% interval at the smaller size. The numbers are CPM's complete gain over best-case DeFi and over
the bank arrangement at both ends, its gain over the pool without the charge at face value, real cost saved against
best-case DeFi, the difference in failed payments against the bank arrangement, owners' return less their
alternative in years 4 and 5, the bridge cost to the borrower, the instalment coupon, $\theta$ and $c$. Read as
written, no size up to 4,000 passes and the rule takes 8,000. The population is set at 4,000 people, a declared
departure: the rule makes 36 comparisons at each doubling with no allowance for their number, and what still moves
at 4,000 moves when a single large loan is written down. The check was repeated after card limits and the rule for
an extra layoff's wage (Section~\ref{sec:model}) were built into the model, with its reading written before the run. Between 4,000 and 8,000 no gain over the bank
arrangement or best-case DeFi, at either end, moves in any of the three cells, nor does real cost or the gap in
failed payments. Four of the 36 numbers move: the charge's edge over the pool without it at $a_{0} = 0.05$ and 0.6,
owners' return at $a_{0} = 0.05$, and $c$ at $a_{0} = 0.6$.

\emph{Seeds for dollar readings.} The seed-count rule asks for the fewest seeds $n \ge 2$ at which the paired $t$ test at
the arrangement's Bonferroni level $\alpha$ shows a true difference of zero clear of a margin $\Delta$ with
probability 0.8:
\begin{equation}
  n^{*} = \min\Bigl\{\, n \ge 2 \;:\; \Pr\bigl[\,T_{n-1}(\lambda_{n}) > t_{1-\alpha/2,\,n-1}\,\bigr] \ge 0.8 \Bigr\},
  \qquad \lambda_{n} = \frac{\Delta\sqrt{n}}{s},
  \label{eq:seeds}
\end{equation}
where $T_{n-1}(\lambda)$ is a noncentral $t$ variable with $n-1$ degrees of freedom. The spread is pooled across the
cells $k$ of a reading, each with $n_{k}$ seeds and spread $s_{k}$,
\begin{equation}
  s = \Bigl(\sum_{k}(n_{k}-1)\,s_{k}^{2} \Big/ \sum_{k}(n_{k}-1)\Bigr)^{1/2}.
  \label{eq:pooled}
\end{equation}
The readings are, for CPM and indexed CPM, the real cost saved against best-case DeFi and against the bank
arrangement, pooled over the level-0 cells with $\Delta = \$2$, and the dip against each, pooled over the cells of
each stressed level with $\Delta$ one failed payment a year, at $\alpha = 0.05/21$ for CPM and $0.05/14$ for indexed
CPM. The largest count is taken, at most 100. A cell whose own spread needs more seeds than it ran is underpowered.

\emph{Yes-or-no readings.} A share of $k$ of $n$ seeds is read by the lower end of its Wilson interval,
\begin{equation}
  L = \frac{\hat{p} + \dfrac{z^{2}}{2n} - z\sqrt{\dfrac{\hat{p}(1-\hat{p})}{n} + \dfrac{z^{2}}{4n^{2}}}}
           {1 + \dfrac{z^{2}}{n}}, \qquad \hat{p} = \frac{k}{n}, \quad z = \Phi^{-1}(1-\alpha/2),
  \label{eq:wilson}
\end{equation}
and the share reads above one half when $L > 1/2$. At $\alpha = 0.05$, $z = 1.96$; under the Bonferroni correction
over 21 cells, $z = 3.04$, and over 14, $z = 2.91$. Twenty seeds need 15 such seeds at 95\% and 17 under the
correction over 21 cells; \SeedsConfirm{} seeds need \WilsonConfirm{} at 95\%, \WilsonConfirmBonf{} under the
correction over 21 cells and \WilsonConfirmBonfIndexed{} over 14; \SeedsExtended{} seeds need \WilsonExtended{},
\WilsonExtendedBonf{} and \WilsonExtendedBonfIndexed{}.

\emph{The confirmation count.} The seed-count rule reads each spread from the development map, so the spread is
itself an estimate. That map has three seeds a cell, so each cell adds only two degrees of freedom to the estimate.
The confirmation count allows for this: each
reading's pooled spread~\eqref{eq:pooled} is taken at the upper end of its one-sided interval at assurance
$q = \SeedAssurance{}$, the level of the power,
\begin{equation}
  s^{+} = s\,\sqrt{\frac{\nu}{\chi^{2}_{1-q}(\nu)}}, \qquad \nu = \sum_{k}(n_{k}-1),
  \label{eq:spreadbound}
\end{equation}
where $\chi^{2}_{1-q}(\nu)$ is the $1-q$ quantile of the chi-square distribution with $\nu$ degrees of freedom, so
the true spread lies below $s^{+}$ with probability $q$. The count is the largest $n^{*}$ of~\eqref{eq:seeds}
at $s^{+}$ over the readings, at most 100. On the final development map the seed-count rule's own count, at the spreads
themselves, is \SeedsPlanned{}, set by the dip against best-case DeFi at 2 points of distress. That reading's spread is
\SpreadBinding{} failed payments a year on $\nu = \FreedomBinding{}$ degrees of freedom, so $s^{+} =
\SpreadBindingBound{}$, and the allowance takes the count to \SeedsAssured{}. The confirmation count is
$n = \SeedsConfirm{}$, fixed by this rule before any confirmation seed was run or read. No cell with no distress
needs more than \SeedsCellMost{} seeds on its own spread, so \SeedsConfirm{} seeds leave \CellsUnderpowered{} cells
underpowered on the development spreads. The \SeedsConfirm{} seeds run one seed a part, each part written when
it is done. Between parts only integrity is checked: that no run failed, that every run of the seed is there and
that money is conserved. The parts are joined into the stored results that a single pass over all the seeds would
write, which a test reproduces value for value. The claim is read once, at \SeedsConfirm{}.
One extension is declared in advance: the cells the seed-count rule marks underpowered at \SeedsConfirm{} may be run to
\SeedsExtended{} seeds, all such cells together, and are then read once at \SeedsExtended{}. On the development
check run after card limits were built into the model, the pooled spread of the cost reading was \$0.99 without stress and \$1.42 at 2
points, against the \SpreadResolved{} up to which \SeedsConfirm{} seeds resolve the \$2 margin at the Bonferroni
level.

\subsection{Development and confirmation}

Three development seeds were used to build and tune the model and to read every development map. The final
development map runs on them at 4,000 people over all 105 cells, and the bounds from development seeds are read from
it. Three earlier blocks of seeds are retired and never read as confirmation: one because review scripts rebuilt
a population from it, one because a confirmation attempt crashed on it and was debugged there with no outcome
stored or read, and the eight seeds of the size check. The confirmation seeds are the first $n$ of a block never run,
and the simulation code refuses a confirmation run until the count is fixed at the freeze. The plan, the count and the code's
commit are frozen together. The confirmation map is then run under the decided entry rule, the quote rule
(entrants read the published quote), and under the record rule (they read the pool's realised twelve-month
return), the pessimistic one. Its outcomes are read once.
Nothing in the code, the plan or the reading changes after a confirmation seed is read, and any later analysis is
labelled exploratory. The headline results of Section~\ref{sec:results} and the registered claim are that
reading. What was run on development seeds is reported in Section~\ref{sec:robustness} and labelled so.
